\section{Source audit, verification, and further research}
\label{qspec:sec:audit}

\subsection{What was inspected}
The repository anchor is \qpin. Its commit timestamp is
11 October 2026 at 05:06:38 UTC, which is 10 October 2026 in the user's
America/Los\_Angeles time zone. The local date in this delivery is therefore
10 October 2026.

The source reading included the manuscript's main file and preface, its
research-status chapter, the beginning and central conventions of the
integration chapter, the displayed $S_6$ target and formal-obstruction
framework in the full-Cayley continuation, the arithmetic spectral-jets
section on periodic weights and poles, and the all-rational formal
classification section of the Herglotz report. The intake README and
incoming directory catalogue were also read \cite{repo,arithmetic,classification,intake}.
This was a targeted reading, not a full audit of every included manuscript
chapter, report, proof prerequisite, or source program.

Six incoming archive names were visible at the pin: the Bernoulli--Tornheim--
digamma, cubic-tail-descent, determinantal-harmonic-resonance,
differential-polygamma-and-Lerch-jets, multishift-reflection, and
polygamma-tail-identities packages. Only their catalogue metadata was
inspected; their interiors were not read. Accordingly, no theorem,
conjecture, omission, or error is attributed here to any of those archive
interiors. The machine-readable provenance file and source-audit note
preserve this distinction.

\subsection{A concrete editorial synchronization correction}
The inspected research chapter still describes the general rational
$J(p/q)$ symbol classification as a further question. At the same pin,
the report \qfile{herglotz-rational-classification} already contains a
section entitled ``Complete formal classification at rational arguments''
with a theorem and its argument \cite{classification}.

Its statement concerns coprime positive $p,q$. In mixed parity, with
$e$ the even entry and $o$ the odd entry, formal rational-dilogarithmic
reducibility and formal logarithmic reducibility are equivalent to
\begin{equation}\label{qspec:eq:existing-J-status}
e^2\equiv\pm1\pmod o,\qquad o^2\equiv1\pmod{2e}.
\end{equation}
When both entries are odd, the stated rational-dilogarithmic cases have
$p,q\in\{1,3,5\}$, and the only logarithmic case is $p=q=1$.
This classification belongs to the existing report; it is not a result
of the present article. Its prerequisite symbol-theoretic lemmas were
not all re-audited here.

The proposed correction is therefore editorial: replace the stale open
classification question by a pointer to that existing theorem, subject
to the repository's normal integration review, while retaining the
separate tasks of constructing short explicit evaluations and understanding
numerical period relations. Failure of a formal reduction criterion is
not numerical transcendence or period independence. The delivered
\qfile{integration/editorial_update.md} contains precise replacement wording.

No new counterexample to an analytic formula in the inspected source
material was established in this session. The following mathematical
safeguards concern the present derivations; they are not allegations
about unread reports.

\subsection{Five distinctions required for a correct implementation}
\begin{enumerate}
\item \textbf{The doubled spectral variable.} The pole of
$\zeta(1+2\epsilon,a)$ is $1/(2\epsilon)$, and its $m$th regular
jet is $(-2)^m\gamma_m(a)$. Omitting the factor two breaks both
reflection channels and the anomaly residue calculation.
\item \textbf{A finite part is not a Taylor coefficient before pole removal.}
The coefficient $C_m$ is
$\Bcal_{m+1}(\ell)/(m+1)$, not $\Bcal_m(\ell)$. The extra index comes
from dividing the analytic function $\epsilon C(\epsilon)$ by
$\epsilon$.
\item \textbf{The binomial radius is real.} Formula
\eqref{qspec:eq:harmonic-counterterms} cannot simply be evaluated outside
$b<\min(a,1-a)$. The finite-head formula supplies the exact replacement.
Both tail shifts increase by $M$.
\item \textbf{Complex-order holomorphy is not a real-part operation.}
The combinations in \eqref{qspec:eq:zt} are needed before differentiating
in a complex spectral order. At real orders they may then be evaluated
as real and imaginary parts.
\item \textbf{Pointwise logarithms do not automatically factor continued sums.}
The odd-weight examples \eqref{qspec:eq:anomaly-examples} are explicit
counterexamples to unchecked multiplicativity. Evenness guarantees
universal vanishing for this conjugate-pair polynomial problem; a
fortuitous zero at one mass does not imply evenness.
\end{enumerate}

\subsection{Executed verification and its evidentiary meaning}
The supplied code was copied to a separate working directory before the
recorded runs. The results are included in
\qfile{data/exact_results.json} and \qfile{data/numeric_results.json}.
The exact suite ran with SymPy 1.14.0. The numerical suite ran with
mpmath 1.3.0 at 45 working decimal digits.

\begin{center}
\begin{tabular}{@{}p{0.65\linewidth}r@{}}
\toprule
Exact finite regression family & Assertions\\
\midrule
Pochhammer jets and Bell polynomials & 140\\
Ordinary/half-shifted harmonic conversion & 120\\
Anomaly coefficients, their nonvanishing, and parity zeros & 36\\
Rational partial fractions & 6\\
Half-integer Bessel recurrence & 12\\
Mass-primitive coefficient recurrence & 85\\
Shift/mass differential equation & 88\\
Heat-residue normalization & 13\\
Anchored primitive pole cancellation & 8\\
Several-mass polynomial change of basis & 9\\
\midrule
Total & 517\\
\bottomrule
\end{tabular}
\end{center}

The numerical suite passed 82 independent comparisons. Eighteen test the
two Stieltjes channels through order three, including masses outside the
unshifted binomial disk. Twenty-four compare negative-point derivatives
using the differentiated Hurwitz series, finite complex-Hurwitz formula,
and finite polylogarithm formula. Sixteen test even and odd polynomial
anomalies; ten compare positive-integer polygamma and polylogarithm
expressions. The remaining fourteen test general Poisson continuation,
Gamma ratios and integrals, parameter primitives, both polylogarithm
integrals, and the several-mass theorem, including repeated masses.

The largest recorded scaled residual was
$4.618\times10^{-44}$, where the scale is
$\max(1,|\mathrm{lhs}|,|\mathrm{rhs}|)$. The acceptance threshold was
$10^{-33}$. These are \emph{floating-point diagnostics}, not rigorous
interval enclosures, tail certificates, or proofs of equality. The
truncation choices are conservative numerical choices, not a delivered
validated error estimator. The exact finite suite proves its particular
algebraic assertions, not the analytic continuation or all-index theorems.
Those depend on the ordinary proofs given in the article.

\subsection{Integration recommendation and claim boundaries}
A suitable destination is a new report under
\qfile{Analysis/Polylogarithms/docs/reports/quadratic-spectrum-stieltjes/}.
The compact integration fragment uses prefixed labels and can be adapted
into the Hurwitz--Stieltjes and integration chapters after normal editorial
review. No automatic patch, remote commit, overwrite of existing research,
or change of conjecture status is included.

The key claims to preserve are: all-index exact identities in both
reflection channels; all-positive-mass convergence through a finite head;
a universal, not pointwise, parity characterization; finite first-derivative
reductions, not arbitrary higher-derivative factorization; and explicitly
anchored primitives. The classical inputs and the new derivations should
remain distinguishable in the merged text.

\subsection{Further research questions and concrete next targets}
\begin{enumerate}
\item \textbf{Higher spectral jets at negative integers.}
For $r\ge2$, determine the smallest explicitly stated function algebra
containing $\partial_s^rE(-m;a,b)$. The reciprocal-Gamma zero leaves
Bessel order derivatives of order at most $r-1$. Can the second jet be
reduced to a finite family of polylogarithm order derivatives plus
explicit logarithmic terms, or is an additional integral coordinate
required? A failure of a chosen formal calculus must be distinguished
from numerical nonreducibility.

\item \textbf{Higher-jet factorization defects.}
For a fixed polynomial $W$, study
\[
\Delta_r(W)=2^{r-1}\bigl(F_{W,+}^{(r)}(0)+F_{W,-}^{(r)}(0)\bigr)
-Z_W^{(r)}(0).
\]
At $r=1$ Theorem~\ref{qspec:thm:anomaly} gives an exact local polynomial.
At $r=2$ mixed logarithms introduce a genuinely nonlocal part. Separate
that part from the residue polynomial and determine what parity still
forces to vanish. This is an extension problem, not an assumption that
the first-derivative criterion persists unchanged.

\item \textbf{Noninteger polynomial degree.}
Replace $W(x)$ by $x^\rho$ on the positive spectrum. Determine the exact
continuation as $\rho$ crosses the resonant odd integers and formulate a
parameter-dependent correction that specializes to
\eqref{qspec:eq:anomaly-harmonic}. The finite polynomial residue in this
article provides a set of exact interpolation constraints.

\item \textbf{Arithmetic projection of both Stieltjes channels.}
At $a=p/q$, apply finite character projectors to the master identities
and their mass-transfer laws. Which combinations admit a shorter
Gamma/Dirichlet-$L$ jet description, and which retain genuinely useful
damped polylogarithm integrals? The first useful target is a complete
small-denominator dictionary with the even and odd character components
kept separate.

\item \textbf{Exact removal of the auxiliary mass.}
The identities are independent of $b$ after their counterterms are
combined. Develop an exact coefficient calculus for the cancellations
as $b\downarrow0$, rather than simply taking a numerically unstable
limit of separate terms. Comparing the resulting coefficients at two
masses may generate new harmonic sums involving Hurwitz order derivatives.

\item \textbf{Mixed shift and mass primitives.}
Combine \eqref{qspec:eq:all-primitives},
\eqref{qspec:eq:mass-primitives}, and \eqref{qspec:eq:PDE} to classify
which rectangular integrals in $(a,b^2)$ have finite polylogarithm and
Hurwitz-jet endpoints. Anchoring every primitive should be part of the
statement, not a later adjustment by unspecified polynomials.

\item \textbf{Several factors beyond the first spectral derivative.}
Theorem~\ref{qspec:thm:several-masses} has a finite answer even with
repeated masses. Determine the second derivative for two masses and
its confluent limit as they coincide. The target is an exact formula
for the logarithmic cross term, with no subtraction of divergent sums
left implicit.

\item \textbf{Formalization and validated evaluation.}
The Bell recurrences, harmonic coefficient identity, rational partial
fractions, and finite mass/PDE reductions are suitable finite targets
for a proof assistant. Separately, an interval implementation of the
Bessel order integrals and finite-head remainder would turn the present
floating-point diagnostics into validated evaluations. Neither has
been supplied here.
\end{enumerate}

\section*{Conclusion}
\addcontentsline{toc}{section}{Conclusion}
The positive quadratic deformation does more than provide another series
for a Stieltjes constant. Its even and odd spectra form a complete
reflection pair, its harmonic counterterms admit a globally valid
finite-head formulation, and its negative spectral zeros collapse first
derivatives to finite polylogarithm expressions. Polynomial parity
explains exactly when conjugate factorization is legitimate. The same
finite family supports anchored mass primitives, a shift/mass differential
equation, and products of several quadratic factors. These are exact
structural identities with ordinary proofs; the supplied computations
check their implementation without replacing those proofs.
